%%%%%%%%%%  scan idx 49  |  folio 36  %%%%%%%%%%
\SGAPrintPage{36}%
ment of $K$ such that $\lambda^iN=0$ for sufficiently large $i$. We can then, for every $x\in K$ and every integer $p\geq1$, consider the elements $\lambda^p(N,x)$, which are expressed by universal polynomials with integer coefficients in the $\lambda^iN$ and the $\lambda^jx$. We have $\lambda^1(N,x)=x$ and we have formula (1.32) \emph{in} $K$. To express the properties of the operations $x\mapsto\lambda^p(N,x)$, it is convenient to introduce a new $\lambda$-ring denoted $K_N$. As a commutative ring, $K_N$ is obtained by adjoining a unit to the ring whose underlying additive group is $K$ and whose multiplication is given by
\sgaStableTarget{sga6.e00.eq.1-34}\begin{equation}\tag{1.34}
x_Ny=xy\mu\qquad (\mu=\lambda_{-1}(N)).
\end{equation}
The map $\lambda_t$ from $K_N$ to $K_N[[t]]$ is defined by:
\sgaStableTarget{sga6.e00.eq.1-35}\begin{equation}\tag{1.35}
\lambda_t(n,x)=(1+t)^n\sum_{p\geq0}\lambda^p(N,x)t^p,
\end{equation}
where we have put $\lambda^0(N,x)=(1,0)$ (the unit element of $K_N$). I say that, \emph{equipped with these operations, $K_N$ is a special $\lambda$-ring}. We therefore have to verify the formulas $\lambda_t(X+Y)=\lambda_t(X)\lambda_t(Y)$, $\lambda_t(XY)=\lambda_t(X).\lambda_t(Y)$, and $\lambda_t(\lambda^i(X))=\lambda^i(\lambda_t(X))$ for $X,Y\in K_N$. We can restrict ourselves to verifying these formulas when $X$ and $Y$ are in $K$, and we find that it suffices to verify them when $K$ is the free special $\lambda$-ring generated by $X,Y$ and $N$ subject to the relations $\lambda^iN=0$ for $i>q$. But the group homomorphism
\sgaStableTarget{sga6.e00.eq.1-36}\begin{equation}\tag{1.36}
K_N\longrightarrow K,
\end{equation}
which takes unit to unit and on $K$ reduces to the homomorphism $x\longrightarrow x\mu$, is a homomorphism compatible with the ring structures and the maps $\lambda^i$, as we verify immediately. In the present case,
%%%%%%%%%%  scan idx 50  |  folio 37  %%%%%%%%%%
\SGAPrintPage{37}%
its
restriction to $K$ is \emph{injective}; as $K$ is a special $\lambda$-ring, it follows readily that the above relations do hold for $X,Y\in K$, and this completes the proof of our assertion.

According to formula (1.23), it is natural, for $x\in K$ and $n\geq1$, to put
\sgaStableTarget{sga6.e00.eq.1-37}\begin{equation}\tag{1.37}
\gamma^n(N,x)=\sum_{i=0}^{n-1}\binom{n-1}{i}\lambda^{n-i}(N,x),
\end{equation}
in other words, we put $\gamma^n(N,x)=\gamma^n(x_N)$, where $x_N$ denotes $x$ considered as an element of the $\lambda$-ring $K_N$.

\sgaStableTarget{sga6.e00.prop.1-5}\textbf{Proposition 1.5.} \emph{Let $A$ be a commutative graded ring with unit, $N$ an element of $\widetilde A$ of the form $N=[q,1+N^1+\cdots+N^q]$, and $x$ an arbitrary element of $\widetilde A$: $x=[\nu,1+\sum_{i\geq1}x^i]$. Then, for every integer $j\geq1$, we have $\gamma^{q+j}(N,x)\in\widetilde A^j$, and the component of degree $j$ of $\gamma^{q+j}(N,x)$ is of the form}
\sgaStableTarget{sga6.e00.eq.1-38}\begin{equation}\tag{1.38}
\left(\gamma^{q+j}(N,x)\right)^{(j)}=(-1)^{j-1}(j-1)!G_{q,j}(\nu,x^1,\ldots,x^j;N^1,\ldots,N^j),
\end{equation}
\emph{where $G_{q,j}$ is a universal polynomial with integer coefficients characterized by the condition}
\sgaStableTarget{sga6.e00.eq.1-39}\begin{equation}\tag{1.39}
G_{q,j}(\nu,x^1,\ldots,N^j)N^q=(-1)^q(x*\lambda_{-1}(N))^{(q+j)}.
\end{equation}

\textbf{Proof.} We can assume that $A$ is the ring of polynomials with integer coefficients in indeterminates $x^i$ and $N^j$ (with $1\leq j\leq q$). As we evidently have, combining (1.32) and (1.37),
\sgaStableTarget{sga6.e00.eq.1-40}\begin{equation}\tag{1.40}
\gamma^n(N,x)\lambda_{-1}(N)=\gamma^n(x\lambda_{-1}(N)),
\end{equation}
and, for $n=q+j$, the second member has filtration $\geq q+j$ (use formula (1.22 bis), taking into account that $x\lambda_{-1}(N)$ has
%%%%%%%%%%  scan idx 51  |  folio 38  %%%%%%%%%%
\SGAPrintPage{38}%
zero augmentation), whereas $\lambda_{-1}(N)$ has filtration $q$ (formula (1.22)), it follows that $\gamma^n(N,x)$ has filtration $j$, since, by (1.18), the associated graded ring of $A$ is here an \emph{integral domain}. Using the preceding formulas, we immediately obtain the indicated form for $(\gamma^{q+j}(N,x))^{(j)}$.
\begin{flushright}q.e.d.\end{flushright}

\sgaStableTarget{sga6.e00.chapter-ii-classes-of-coherent-algebraic-sheaves.occ2}\section*{CHAPTER II}
\begin{center}
Classes of coherent algebraic sheaves\par
and Chern classes
\end{center}

\sgaStableTarget{sga6.e00.sec.1.occ3}\subsection*{\S 1. Chow's theory}

Throughout this chapter, we assume a \emph{base field} $k$ fixed, which we shall assume algebraically closed, for simplicity. A large part of what follows, and perhaps all of it, is nevertheless valid without this restriction. The algebraic spaces and regular maps considered will be defined over $k$.

An algebraic space is called \emph{quasi-projective} if it is isomorphic to a locally closed subset of a projective space.

Let $X$ be a connected nonsingular quasi-projective variety of dimension $n$. We denote by $A(X)$, and call the \emph{Chow ring} of $X$, the graded ring of classes of cycles on $X$ modulo rational equivalence. ($A^i(X)$ consists of the classes of cycles of dimension $n-i$). The fact that $A(X)$
%%%%%%%%%%  scan idx 52  |  folio 39  %%%%%%%%%%
\SGAPrintPage{39}%
is indeed a graded ring follows from the following theorem of Chow:

\sgaStableTarget{sga6.e00.thm.2-1}\textbf{Theorem 2.1.}\SGAAuthorFootnote{*}{As a result of a misunderstanding on the author's part, the theorem states more than is currently known. In fact, only Chow's ``moving lemma'' is used, saying that two cycles can be moved within their classes so that their intersection has no excess; cf.\ Séminaire Chevalley 1958, Anneau de Chow et applications.} \emph{Every cycle on $X$ is rationally equivalent to a nonsingular cycle. Given elements of $A^i(X)$ and $A^j(X)$, we can find nonsingular representatives of these elements, say $Y^{n-i}$ and $Z^{n-j}$, which meet transversely everywhere.}

(A cycle is called nonsingular if its components are nonsingular subvarieties. We say that two cycles meet transversely everywhere if every component of one meets every component of the other transversely everywhere, i.e. if the two intersecting varieties are nonsingular at their points of intersection and their tangent varieties there are in general position \dots)

Let $f$ be a morphism (i.e. a regular map) from $X$ to $Y$ ($X$ and $Y$ being quasi-projective and nonsingular); then $f$ defines (using the preceding theorem) a homomorphism of graded rings:
\sgaStableTarget{sga6.e00.eq.2-1.occ2}\begin{equation}\tag{2.1}
f^*:A(Y)\longrightarrow A(X)
\end{equation}
(\emph{inverse image of cycle classes}). Thus $A(X)$ becomes a contravariant functor in $X$.

Let $f$ be a \emph{proper} morphism from $X^n$ to $Y^m$ (i.e. in Weil's terminology, the graph of $f$ is complete over $Y$, cf.\ [3]). Then $f$ also defines an additive homomorphism
\sgaStableTarget{sga6.e00.eq.2-2}\begin{equation}\tag{2.2}
f_*:A(X^n)\longrightarrow A(Y^m)
\end{equation}
preserving the dimension of cycles, and therefore increasing the degree by $m-n$.

% --- Source pages 47--75 ---
%%%%%%%%%%  scan idx 53  |  folio 40  %%%%%%%%%%
\SGAPrintPage{40}%
Thus, with respect to proper morphisms, and disregarding the multiplicative structures and gradings, $A(X)$ is a covariant functor in $X$. Let us note the classical formula
\sgaStableTarget{sga6.e00.eq.2-3.occ2}\begin{equation}\tag{2.3}
 f_*(x.f^*(y))=f_*(x).y .
\end{equation}

We shall also use Chow's theory of Chern classes. It consists in attaching to every vector bundle $E$ on $X$ (still assumed quasi-projective and nonsingular) \emph{Chern classes} $C^i(E)\in A^i(X)$, ($i\geq 1$), in such a way as to satisfy the following conditions; put $C^0(E)=1$ and
\sgaStableTarget{sga6.e00.eq.2-4.occ2}\begin{equation}\tag{2.4}
 \Ctil(E)=\left[\rang E,\sum_{i\geq 0}C^i(E)\right]\in \Atil(X)
\end{equation}
(see chapter I, \S 3 for the definition of $\Atil(X)$). The conditions to be satisfied by a theory of Chern classes are then expressed as follows:
\sgaStableTarget{sga6.e00.eq.2-5.occ2}\begin{equation}\tag{2.5}
 \Ctil(E)=\Ctil(E')+\Ctil(E'')
\end{equation}
if the vector bundle $E$ is an extension of $E'$ by $E''$ (putting $C(E)=\sum_{i\geq0}C^i(E)$, this condition is expressed by $C(E)=C(E').C(E'')$). Formula (2.5) allows $\Ctil$ to be extended to an additive homomorphism from $\underline{K}(\xi(X))$ (cf.\ Chap.\ I, \S 2, example 2) to $\Atil(X)$:
\sgaStableTarget{sga6.e00.eq.2-6}\begin{equation}\tag{2.6}
 \Ctil:\underline{K}(\xi(X))\longrightarrow \Atil(X),
\end{equation}
and we want $\Ctil$ to be a \emph{$\lambda$-homomorphism of $\lambda$-rings}; this last condition is also written
\sgaStableTarget{sga6.e00.eq.2-7}\begin{equation}\tag{2.7}
 \Ctil(E\otimes F)=\Ctil(E)\Ctil(F)\qquad
 \Ctil(\Lambda^iE)=\lambda^i\Ctil(E),
\end{equation}
for two vector bundles $E$ and $F$ (of course, $\Ctil$ is automatically compatible with the augmentations). Let $D$ be a divisor on $X$ and $L(D)$ the vector bundle it defines; we want
